\documentclass[10pt]{article}

\usepackage[utf8]{inputenc}

\usepackage[T1]{fontenc}

\usepackage{amsmath}

\usepackage{amsfonts}

\usepackage{amssymb}

\usepackage[version=4]{mhchem}

\usepackage{stmaryrd}

\usepackage{mathrsfs}

\usepackage{bbold}

\usepackage{graphicx}

\usepackage[export]{adjustbox}

\graphicspath{ {./images/} }



\begin{document}

ми $H_{c}^{q}(X, \mathscr{F})$ конечны и равны 0 шри $q>2 \operatorname{dim} X$. Если к тому же $X$-аффинное многообразие, то $H^{q}(X, \mathscr{F})=0$ для $q>\operatorname{dim} X$.



Для многообразий над полем комплексных чисел Ә. к. конструктивных пучков совпадают с-классич. когомологиями со значениями в этих пучках. Справедлива теорем а о специали зации для г ла дк о го мо р ф и м а: пусть $f: X \rightarrow Y$ - гладкий собственный морфизм схем, и целое число $n$ обрати«о на $Y$; тогда пучки $R^{q} f_{*}(\mathbb{Z} / n \mathbb{Z})$ локально постоянны на $Y$.



Для Э. к. имеют место аналог двойственности Пуанкаре (см. Двойственность в алгебраической геометрии) и Кюннета бормуль. Каждый алгебраич. цикл коразмерности $i$ дает класс когомологий в размерности $2 i$, что позволяет построить теорию Чжәня классов.



Э. к. конструктивных пучков используются для построения l-aдических когомологий и доказательства гипотез Вейля о дзета-функции.



Лит.: [1] Г р о т е н д и к А., в кн.: Международный математический конгресс в Әдинбурге. Сб. докладов, М., 1962, с. 116-37; [2] м и л н Д ж., Этальные когомологии, пер. с aнгл., M., 1983; [3] Cohomologie étale B.- [e.a.], 1977; [4] des topos et cohomologie étale des schémas, t. $1-3$, B. - [e.a] 1972-73. B. И. Данилов



ӘТАЛЬНЫЙ МОРФИЗМ - гладкий морфизм алгебраич. многообразий или схем относительной размерности 0. Эквивалентным образом можно определить Ә. м. схем $f: X \rightarrow Y$ как локально конечно представленный шлоский морфизм такой, что для любой точки $y \in Y$ $k(y)$-схема $\left.f^{-1}(y)=X \bigotimes\right)^{k}(y)$ конечна и сепарабельна. Э. м. обладает свойством подъема инфинитезимальных деформаций: если $f: X \rightarrow Y-$ М. м., $Y^{\prime}-$ аффинная $Y$-схема и $Y_{0}^{\prime}$ - замкнутая подсхема в $Y^{\prime}$, задаваемая нильпотентным пучком идеалов, то естественное отображение $\mathrm{Hom}_{Y}\left(Y^{\prime}, X\right) \rightarrow \operatorname{Hom}_{Y}\left(Y_{0}^{\prime}, X\right)$ биективно. Указанное свойство характеризует Э. м. Наконец, Э. м. можно определить как плоский неразветвленный морфизм. (Локально конечно представленный морфизм $f: X \rightarrow Y$ н е р а 3 в е т в л е н, если диагональное вложение $X \rightarrow X \times{ }_{Y} X$ является локальным изоморфизмом.)



Әтальность (как и гладкость и неразветвленность) сохраняется при композиции морфизмов и при замене базы. Открытое вложение является Ә. м. Любой морфизм между этальными $Y$-схемами этальный. Для гладких многообразий этальность $f: X \rightarrow Y$ означает, что $f$ индуцирует изоморфизм касательных пространств. Локально Э. м. задается многочленом с ненулевой производной.



Э. м. играют важную роль в теории этальных когомологий, при определении фундалентальной әруппь схемы, алгебрацческого пространства и гензелева кольца.



Jum.: [1] G r ot he $n$ d i e c k A., D i e u d on n e J., Eléments de géométrie algébrique, N. Y., 1971; [2] Revêtements étales et groupe fondamental, B.- [e. a.], 1971. В. И. Данилов.



ӘФФЕКТИВНАЯ ОЦЕНКА - несмещенная статистическая оценка, дисперсия к-рой совпадает с нижней гранью в Рао - Крамера неравенстве. Э. о. является достаточной статистикой для оцениваемого параметра. Если Э. о. существует, то ее можно получить с помощью метода максимального правдоподобия. В силу того, что во многих случаях нижняя грань в неравенстве Рао - Крамера не является достижимой, в математич. статистике часто Ә. о. наз. оценку, имеющую минимальную дисперсию в классе всех несмеценных оценок рассматриваемого параметра.



Jum.: [1] K р а м е p Г., Математические методы статистики, пер. с англ., 2 изд., М., 1975; [2] И 6 р а г и мо в И. А., X а с ь м ин с к и й Р.' 3., Асимптотическая теория оценивания, М., 1979; [3] Р а о С. Р., Линейные статистические методы и их применения, пер. с англ., М., 1968. М. С. Никулин.



ӘФФЕКТИВНОСТЬ АСИМПТОТИЧЕСКАЯ КРИТЕРИЯ - понятие, позволяющее осуществлять в случае больших выборок количественное сравнение двух различных статистич. критериев, применяемых для проверки ложной и той же статистич. гипотезы. Необходимость измерять эффективность критериев возникла в 30-40-е гг., когда появились простые с точки зрения вычислений, но «неəфффегтивные» ранговые процедуры.



Существует несколько различных подходов к определению асимптотич. әффективности критерия. Пусть распределение наблюдений определяется действитель ным параметром $\theta$ и требуется проверить гипотезу $H_{0}$ : $\theta=\theta_{0}$ против альтернативы $H_{1}: \theta \neq \theta_{0}$. Пусть также нек-рому критерию с уровнем значимости $\alpha$ требуется $N_{1}$ наблюдений для достижения мощности $\beta$ против заданной альтернативы $\theta$, а другому критерию того же уровня требуется для этого $N_{2}$ наблюдений. Тогда можно определить о т н о и т е л ь н у ю әффф кт и в н о с т б первого критерия по отношению ко второму формулой $e_{12}=\frac{N_{2}}{N_{1}}$. Понятие относительной әффективности дает исчерпывающую информацию при сравнении критериев, но оказывается неудобным для применения, поскольку величина $e_{12}$ является функцией трех аргументов: $\alpha, \beta$ и $\theta$ и обычно не поддается вычислению в явном виде. Для того чтобы обойти әту трудность, используют предельные переходы.



Величина $\lim _{\theta \rightarrow \theta_{0}} e_{12}(\alpha, \beta, \theta)$ при фикспрованных $\alpha$ и $\beta$ (если этот предел существует) наз. а с и мп т о т и-



\includegraphics[max width=\textwidth, center]{2023_07_09_17b576b714b05d1143a8g-1}

с т ь ю (АОӘ) п о П и т м н у. Аналогично определяется АО по Бахаду ру - этом случае при фиксированных $\beta$ и $\theta$ устремляют к нулю $\alpha$, и А О Ә п о Х о д ж е с у - Л е м а н у, когда при фиксированных $\alpha$ и $\theta$ вычисляют предел при $\beta \rightarrow 1$.



Каждое из әтих определений имеет свои достоинства и недостатки. Так, напр., АОЭ по Питмену вычисляется обычно легче, чем АОЭ по Бахадуру (вычисление последней связано с нетривиальной задачей изучения асимптотики вероятностей больших уклонений тестовых статистик), однако в ряде случаев оказывается менее чувствительным средством для сравнения двух критериев.



Пусть, напр., наблюдения распределены по нормальному закону со средним $\theta$ и дисперсией 1 и проверяется гипотеза $H_{0}: \theta=0$ против альтернативы $H_{1}$ : $\theta>0$. И пусть рассматриваются критерии значимости, основанные на выборочном среднем $\bar{X}$ и дроби Стьюдента $t$. Поскольку $t$-критерий не использует информации о величине дисперсии, он должен проигрывать оптимальному критерию, основанному на $\bar{X}$. Однако с точки зрения АОӘ по Питмену эти критерии эквивалентны. Сдругой стороны, АОЭ ио Бахадуру $t$-критерия по отношению к $\bar{X}$ строго меньше 1 при любом $\theta>0$.



В более сложных случаях АОӘ по Питмену может зависеть от $\alpha$ или $\beta$ и ее вычисление оказывается очень трудным. Тогда вычисляют ее предельное значение при $\beta \rightarrow 1$ или $\alpha \rightarrow 0$. Последнее обычно совпадает с предельным значением АОЭ по Бахадуру при $\theta \rightarrow \theta_{0}$. O других подходах к определению АОӘ см. в [2], [4], последовательные аналоги понятия АОЭ введены в [5] - [7]. Выбор того или иного определения АОӘ должен основываться на том, какая из них дает более точное приближение к относительной эффективности $e_{12}$, однако в настоящее время (1984) в этом направлении потти ничего не известно.



Лит.: [1] К е н д а л л М., С т ь ю а р т А., Статистические выводы и связи, пер. с англ., м., 1973; [2] B a h a d u r R., "Ann. Math. Stat.», 1967, v. 38, № 2, p. 303-24; [3] H o d g e s J.,

Le h m a n n E., там же, 1956, v. 27, № 2, p. 324-35; [4] L e h m a n n E., там же, 1956, v. 27, № 2, р. $324-35 ;$ [4]

P a o C. Р., Јинейные статистические методы и их применения, пер. с англ., М., 1968; [5] L a i T. L., "Ann. Stat.», 1978, 1978, v. 6, ㄴ 3, p. $567-81$; [7] B e r k R., там ке, 1976, v. 4,





\end{document}